\hypertarget{appendix-3b88c6cd7195}{}
\section*{Thread-local storage}
\section{Thread-local storage}\label{thread-local-storage}

\textbf{Curriculum depth:} Introduced · \textbf{First seen:} C007 · \textbf{Latest development:} C007

\subsection{Definition and mechanism}\label{definition-and-mechanism}

thread\_local gives a separate object/storage instance per thread; block scope does not make it automatic. Initialization, access costs and destruction follow implementation/language rules. Thread-local state does not make shared pointees thread-safe, and escaped addresses need a valid thread/object lifetime.

\subsection{Prerequisites}\label{prerequisites}

\begin{itemize}
\tightlist
\item
  \href{}{Storage duration and object lifetime (05-Engine-I, appendix-0b104bcba06b)} --- Thread duration is a distinct backing-storage category, not a synonym for a local name's scope.
\end{itemize}

\subsection{Used by}\label{used-by}

Later chunks extend this same node; dependent concepts are listed here as they are integrated.

\subsection{Recall and next depth}\label{recall-and-next-depth}

Give an example, explain why the rule exists, and identify which guarantee would be missing if the prerequisite were ignored.

\textbf{Pending:} Deeper ownership, pointer/reference, generic programming and module rules from later sources. This is written coverage, not demonstrated personal mastery.

\subsection{Source and navigation}\label{source-and-navigation}

\href{Continuous\%20Notes/Computing\%20Foundations\%20-\%20From\%20Bits\%20to\%20C++\%20and\%20Systems\%20APIs\#A\%20name\%20does\%20not\%20decide\%20how\%20long\%20storage\%20lasts}{C007 teaching chapter} · \hyperref[appendix-fd090d982a15]{Source coverage} · \href{}{Dependency Map (02-Foundations-II, appendix-4ad47bdd269a)}
